We presented an evidence-constrained workflow for diagnostic assistance on unlabeled neutron-current time series. The workflow uses deterministic analysis probes under a persistent workspace contract to produce auditable observations, then turns a structured alert summary into manual-grounded, citation-bearing explanatory support. Its output is deliberately framed as an aid to engineering review rather than an autonomous fault decision.

The retrieval study identifies a practical governance issue: for the evaluated corpus and embedding model, cosine similarity can improve while the priority of pre-designated IAEA sources declines as same-domain supplementary manuals are added. Short Chinese queries are particularly sensitive in the tested setting, and document-level relevance is more forgiving than exact chunk identity under overlapping segmentation. The separate CPU benchmark further shows that multilingual semantic alignment and authoritative-source preference are complementary requirements rather than interchangeable objectives.

Together, these results favor terminology-rich report-derived queries and explicit source policies over similarity-only deployment choices. The study's boundaries, including the absence of a sample-level fault ground truth and limited retrieval and embedding coverage, are stated in \cref{sec:limitations}. Future work will extend the evaluation with expert-reviewed cases, relevance annotations, broader multilingual retrieval configurations, and generation-side factuality assessment.
